\newpage
\begin{center}
\large\textbf{Минобрнауки России}

\large\textbf{Юго-Западный государственный университет}
\vskip 1em
\normalsize{Кафедра программной инженерии}
\vskip 1em
\ifВКР{
        \begin{flushright}
        \begin{tabular}{p{.4\textwidth}}
        \centrow УТВЕРЖДАЮ: \\
        \centrow Заведующий кафедрой \\
        \hrulefill \\
        \setarstrut{\footnotesize}
        \centrow\footnotesize{(подпись, инициалы, фамилия)}\\
        \restorearstrut
        «\underline{\hspace{1cm}}»
        \underline{\hspace{3cm}}
        20\underline{\hspace{1cm}} г.\\
        \end{tabular}
        \end{flushright}
        }\fi
\end{center}
\vspace{1em}
  \begin{center}
  \large
\ifВКР{
ЗАДАНИЕ НА ВЫПУСКНУЮ КВАЛИФИКАЦИОННУЮ РАБОТУ
  ПО ПРОГРАММЕ БАКАЛАВРИАТА}
  \else
ЗАДАНИЕ НА КУРСОВУЮ РАБОТУ (ПРОЕКТ)
\fi
\normalsize
  \end{center}
\vspace{1em}
{\parindent0pt
  Студента \АвторРод, шифр\ \Шифр, группа \Группа
  
1. Тема «\Тема\ \ТемаВтораяСтрока»
\ifВКР{
утверждена приказом ректора ЮЗГУ от \ДатаПриказа\ № \НомерПриказа
}\fi.

2. Срок предоставления работы к защите \СрокПредоставления

3. Исходные данные для создания программной системы:

3.1. Перечень решаемых задач:}

\renewcommand\labelenumi{\theenumi)}

\begin{enumerate}
	\item проанализировать предметную область поиска работы и подбора IT-специалистов;
	\item определить требования к онлайн-платформе и описать основные варианты использования программной системы;
	\item спроектировать архитектуру программной системы, пользовательский интерфейс и структуру базы данных;
	\item реализовать web-платформу для работы с пользователями, профилями соискателей, компаниями, вакансиями, откликами, избранными вакансиями и навыками;
	\item обеспечить разграничение доступа для ролей соискателя, работодателя и администратора;
	\item провести тестирование основных сценариев работы программной системы.
\end{enumerate}


{\parindent0pt
  3.2. Входные данные и требуемые результаты для программы:}

\begin{enumerate}
	\item Входными данными для программной системы являются: регистрационные и авторизационные данные пользователей; сведения о профилях соискателей; данные компаний работодателей; сведения о вакансиях, включая требования, условия работы, статус публикации и связанные навыки; данные откликов соискателей на вакансии; данные об избранных вакансиях; параметры поиска и фильтрации вакансий и профилей соискателей; данные, используемые администратором при управлении пользователями и вакансиями.
	\item Выходными данными для программной системы являются: списки опубликованных вакансий и карточки вакансий; профили соискателей и сведения о компаниях; результаты поиска и фильтрации вакансий и кандидатов; списки откликов и их текущие статусы; списки избранных вакансий пользователя; сведения о вакансиях работодателя и откликах на них; административные списки пользователей и вакансий; сообщения системы о результатах выполненных действий и ошибках при некорректном вводе данных.
\end{enumerate}


{\parindent0pt

4. Содержание работы (по разделам):

4.1. Введение: актуальность темы, цель и задачи работы, краткая характеристика разрабатываемой программной системы.

4.2. Анализ предметной области: характеристика сферы поиска работы и подбора IT-специалистов, роли пользователей, проблемы существующих подходов, обоснование необходимости разработки web-платформы, выделение ключевых сущностей и анализ существующих решений.

4.3. Техническое задание: основание для разработки, цель и назначение программной системы, задачи разработки, требования к данным, функциональные требования, требования к интерфейсу, надёжности, информационной безопасности и программной совместимости, варианты использования.

4.4. Технический проект: архитектура программной системы, обоснование выбора технологий, проектирование клиентской и серверной частей, API-маршрутов, структуры базы данных и порядка развёртывания.

4.5. Рабочий проект: спецификация компонентов программной системы и системное тестирование основных сценариев работы.

4.6. Заключение.

4.7. Список использованных источников.


5. Перечень графического материала:

\списокПлакатов

\vskip 2em
\begin{tabular}{p{6.8cm}C{3.8cm}C{4.8cm}}
Руководитель \ifВКР{ВКР}\else работы (проекта) \fi & \lhrulefill{\fill} & \fillcenter\Руководитель\\
\setarstrut{\footnotesize}
& \footnotesize{(подпись, дата)} & \footnotesize{(инициалы, фамилия)}\\
\restorearstrut
Задание принял к исполнению & \lhrulefill{\fill} & \fillcenter\Автор\\
\setarstrut{\footnotesize}
& \footnotesize{(подпись, дата)} & \footnotesize{(инициалы, фамилия)}\\
\restorearstrut
\end{tabular}
}

\renewcommand\labelenumi{\theenumi.}
